\documentclass{standalone}
\usepackage{tikz}
\usepackage{pgfplots}
\usepackage{pgfplotstable}
\usepgfplotslibrary{groupplots}
\usepackage[utf8]{vietnam}
\pgfplotsset{compat=1.17}
\usepackage{times}

% Định nghĩa hàm gauss(mu, sigma)
\pgfmathdeclarefunction{gauss}{2}{%
	\pgfmathparse{1/(#2 * sqrt(2 * pi)) * exp(-((x - #1)^2) / (2 * #2^2))}%
}

\begin{document}
	
	\begin{tikzpicture}
		\begin{groupplot}[
			group style={group size=2 by 1, horizontal sep=1cm},
			width=6cm, height=6cm,
			domain=-6:6, samples=200,
			anchor=south west,
						title style={at={(0.5,-0.15)}, anchor=north}, % đặt tiêu đề phía dưới
			]
			
			% ----- Frame 1: Trung bình có trọng số -----
			\nextgroupplot[title={(a)}]
			\addplot [color=black!45, line width=1.5pt, dashed] {gauss(-2,1)};
			
			\addplot[color=black!45, line width=1.5pt, dashed] {gauss(0,1)};
			
			\addplot [color=black!45, line width=1.5pt, dashed] {gauss(2,1)};
			
			\addplot [color=black!95, line width=2pt, fill=black!45, opacity=.5] {(gauss(-2,1) + gauss(0,1) + gauss(2,1))/3};
			
			% ----- Frame 2: Hàm Gauss đại diện -----
			\nextgroupplot[title={(b)}]
			\addplot [color=black!45, line width=1.5pt, dashed] {gauss(-2,1)};
			
			\addplot[color=black!45, line width=1.5pt, dashed] {gauss(0,1)};
			
			\addplot [color=black!45, line width=1.5pt, dashed] {gauss(2,1)};
			
			\addplot [color=black!95, line width=2pt, fill=black!45, opacity=.5]{gauss(0,1)};
			
		\end{groupplot}
	\end{tikzpicture}
	
\end{document}
